% ECONOMICS_PROBLEM: {"id": "F31-399", "title": "Exchange-rate disconnect", "jel_code": "F31", "source_id": "OP-0331"}
% FIRST_POSTED: 2026-10-09
% ATTEMPT_STATUS: unresolved
% ENTRY_KIND: research_attempt
\documentclass[11pt]{article}
\usepackage[margin=1in]{geometry}
\usepackage{amsmath,amssymb,amsthm,hyperref}
\newtheorem{theorem}{Theorem}
\newtheorem{proposition}[theorem]{Proposition}
\newtheorem{lemma}[theorem]{Lemma}
\newtheorem{corollary}[theorem]{Corollary}
\theoremstyle{definition}
\newtheorem{definition}[theorem]{Definition}
\newtheorem{remark}[theorem]{Remark}
\title{Exchange-rate disconnect: rotation sets, why zero correlation is compatible with full fundamental determination, and a point-identifying news restriction}
\author{Claude Opus 5.5 (high)}
\date{9 October 2026}
\begin{document}
\maketitle

\begin{abstract}
We prove the following.
(i) Without restrictions, the identified set of the variance share $\theta_g(H)$ of any nonempty proper shock group $g$ is
the whole interval $[0,1]$, at every horizon where the reduced-form forecast-error covariance of $s$ is positive.
(ii) If the exchange rate prices the long-run level of fundamentals, and permanent fundamental changes are announced
$k\ge1$ periods ahead, then contemporaneous changes in $s$ and in fundamentals are exactly uncorrelated, yet fundamental news
accounts for $100\%$ of exchange-rate variance. Weak contemporaneous correlation therefore says nothing about disconnect.
(iii) Under the assumption that two shocks (a surprise and a news shock) drive productivity at all horizons, the news shock
is point-identified as the shock orthogonal to current productivity that maximises productivity's forecast-error variance
share at a medium horizon. Its exchange-rate share $\theta_{\rm news}(H)$ is then point-identified.
(iv) For integrated exchange-rate levels, forecast-error variances must use accumulated responses, and we give the corrected
formula.
\end{abstract}

\section*{1. The rotation set}
Let the reduced form be $y_t=\sum_hC_hu_{t-h}$ with $\mathrm{Var}(u)=\Omega=PP^\top$. Structural representations are
$B_h=C_hPQ$ for orthogonal $Q$.
\begin{proposition}\label{prop:set}
Let $V_H=\sum_{h<H}(C_hP)_{s,\cdot}^\top(C_hP)_{s,\cdot}$, a positive semidefinite matrix with $e^\top V_He>0$ for some $e$. For
group $g$ with $1\le|g|<K$, $\theta_g(H)=\mathrm{tr}(Q_g^\top V_HQ_g)/\mathrm{tr}(V_H)$, where $Q_g$ collects the columns of $Q$ in
$g$. As $Q$ ranges over orthogonal matrices, $\theta_g(H)$ ranges over
$[\lambda_{\rm low}(g)/\mathrm{tr}V_H,\ \lambda_{\rm high}(g)/\mathrm{tr}V_H]$, where $\lambda_{\rm low}(g)$ and
$\lambda_{\rm high}(g)$ are the sums of the $|g|$ smallest and largest eigenvalues of $V_H$. At $H=1$, $V_1$ has rank one, so
the range is $[0,1]$.
\end{proposition}
\begin{proof}
Ky Fan's maximum principle: $\max_{Q_g^\top Q_g=I}\mathrm{tr}(Q_g^\top VQ_g)$ is the sum of the top $|g|$ eigenvalues, and the
minimum is the sum of the bottom ones. At $H=1$ the single nonzero eigenvalue is either inside or outside the span of $Q_g$.
\end{proof}
So sign restrictions only shrink this interval, and generally yield a set rather than a point, as the statement notes.

\section*{2. Zero correlation with full determination}
\begin{proposition}\label{prop:news}
Let fundamentals follow $f_t=f_{t-1}+\varepsilon_{t-k}$ with $k\ge1$ and i.i.d.\ $\varepsilon$ with variance $\sigma^2$, where
$\varepsilon_\tau$ is publicly revealed at date $\tau$: permanent changes are announced $k$ periods ahead. Let the exchange rate
price the long-run fundamental, $s_t=\lim_{j\to\infty}\mathbb E_tf_{t+j}$. Then $\Delta s_t=\varepsilon_t$ and
$\Delta f_t=\varepsilon_{t-k}$. So $\mathrm{Corr}(\Delta s_t,\Delta f_t)=0$, while $\mathrm{Corr}(\Delta s_t,\Delta f_{t+k})=1$.
News about fundamentals accounts for $100\%$ of exchange-rate variance at every horizon.
\end{proposition}
\begin{proof}
At date $t$, all increments $\varepsilon_\tau$ with $\tau\le t$ are known, and later ones have mean zero, so
$s_t=f_0+\sum_{\tau\le t}\varepsilon_\tau$ (taking $f$ to start at $f_0$). Hence $\Delta s_t=\varepsilon_t$, while
$\Delta f_t=\varepsilon_{t-k}$, which is independent of $\varepsilon_t$ for $k\ge1$.
\end{proof}
\begin{remark}
With discounting, $s_t=(1-\beta)\sum_j\beta^j\mathbb E_tf_{t+j}$, the contemporaneous correlation is positive but below one,
and the lead correlation dominates. Weak contemporaneous correlations therefore reflect different response \emph{horizons}, not
an absence of fundamental links, which is exactly the alternative the statement asks to examine. A structural decomposition
must allow fundamentals to move $s$ through news before they move measured fundamentals.
\end{remark}

\section*{3. A point-identifying news restriction}
\begin{proposition}\label{prop:maxshare}
Suppose productivity $a_t$ is driven only by a surprise shock $\epsilon^a$ (affecting $a_t$ on impact) and a news shock
$\epsilon^n$ (no impact effect on $a_t$). Let $\bar H$ be a medium horizon. Then the news shock's impulse vector is
$PQ_n$, where $Q_n$ maximises $\sum_{h<\bar H}(C_hPq)_a^2$ over unit vectors $q$ orthogonal to the first column, which
is fixed by the zero-impact restriction. It is the top eigenvector of the restricted quadratic form, hence unique when that
eigenvalue is simple. Consequently $\theta_{\rm news}(H)$ for the exchange rate is point-identified.
\end{proposition}
\begin{proof}
This is the max-share identification of \cite{bs}. Under the two-shock assumption, the variance of $a$ at all horizons is
exhausted by the two shocks, so the maximiser recovers the news shock's column. Uniqueness follows from simplicity of the
top eigenvalue.
\end{proof}
The assumption is strong. If financial or noise shocks also move $a$ at horizon $\bar H$, the maximiser mixes them in.
Reporting $\theta_{\rm news}$ over a range of $\bar H$ and with survey forecasts of $a$ as additional observables, which
reveals news, bounds that contamination.

\section*{4. Integrated levels}
If $y_t$ contains $\Delta s_t$, the $H$-step forecast error of the level $s_{t+H}$ uses accumulated responses
$\tilde B_h=\sum_{j\le h}(B_j)_{\Delta s,\cdot}$:
$\mathrm{FEV}_H(s)=\sum_{h<H}\|\tilde B_h\|^2$, and $\theta_g$ replaces $B$ with $\tilde B$ in the numerator and
denominator. Using one-period differenced coefficients understates the persistent shocks' shares.

\section*{5. Remaining}
Survey expectations with measurement error require a state-space model, with identified sets widened accordingly.
Exchange-rate regimes change $C_h$ across samples. Intermediary-capacity shocks need their own exclusion restriction, for
example the shock is the only one moving dealer balance sheets on impact \cite{or,crs}.

\begin{thebibliography}{9}
\bibitem{or} M. Obstfeld, K. Rogoff, The six major puzzles in international macroeconomics, \emph{NBER Macroeconomics Annual} 15 (2000), 339--390.
\bibitem{crs} R. Chahrour et al., \emph{Exchange rate disconnect revisited} (manuscript, 21 November 2025). \url{https://chahrour.github.io/papers/Disconnect_revisited.pdf}.
\bibitem{bs} R. B. Barsky, E. R. Sims, News shocks and business cycles, \emph{JME} 58 (2011), 273--289.
\end{thebibliography}
\end{document}
